\SourceLeaf{313}
Da nun die Nenner in II. und III. beide denselben Werth, nemlich

\[
\tfrac{1}{2}\rho\left\{\sin(2A-2A')+\sin(2A'-2A'')+\sin(2A''-2A)\right\},
\]

haben, so ergibt sich aus diesen Gleichungen durch Division:

\[
\operatorname{tg}\psi = -\dfrac{mmh\cos A+m'm'h'\cos A'+m''m''h''\cos A''}{mmh\sin A+m'm'h'\sin A'+m''m''h''\sin A''},
\]

wodurch für $\psi$ und damit auch für $R$ zwei Werthe folgen.]

\bigskip
\begin{center}
\hrulefill
\end{center}
\bigskip

\Needspace{9\baselineskip}
\begin{center}
BEMERKUNG.
\end{center}

Die vorstehende Notiz, die sich in einem Handbuche findet und vermuthlich aus dem Jahre 1816 stammt, ist sehr flüchtig geschrieben, sodass zahlreiche Schreibfehler in den Formeln verbessert werden mussten. Sie bricht bei der Formel II der zweiten Auflösungsart ab, bei der, was im Druck nicht ersichtlich gemacht werden konnte, der Nenner und der Factor $\Delta$ ergänzt werden mussten. Der Deutlichkeit wegen sind dem Texte zwei Figuren hinzugefügt worden.

Es ist anzunehmen, dass \textsc{Gauss} seine Auflösungen gefunden hat, indem er von der Darstellung der Ecken des Dreiecks durch complexe Grössen ausging, denn hierbei ergeben sich genau die Formeln des Textes, wenn man darin die Paare von Gleichungen, die den Cosinus und den Sinus desselben Winkels enthalten, mittelst des Symboles $i$ immer in eine einzige Gleichung zusammenfasst. Dass \textsc{Gauss} die Darstellung der Punkte einer Ebene durch complexe Grössen schon früh bei geometrischen Aufgaben benutzt hat, zeigt zum Beispiel die auf der folgenden Seite abgedruckte Notiz über die Construction der \textsc{Pothenot}schen Aufgabe, sowie eine handschriftliche Bemerkung zu \textsc{Schumachers} Übersetzung der \emph{Geometrie der Stellung} von \textsc{Carnot} aus dem Jahre 1810, die sich ebenfalls auf die Geometrie des Dreiecks bezieht und bereits Bd.\ IV.\ S.\ 396 bis 397 Platz gefunden hat.

\hfill \textsc{Stäckel}.

\bigskip
\begin{center}
\hrulefill
\end{center}
